\section{Discussion}\label{sec:discussion}

\subsection{A causal layer that complements monitoring and machine learning}
The demonstrated capability is not better interpolation of observed concentrations; it is the separation of an observed concentration into mechanistically distinct contributions. On the Rogue, the calibrated model partitions the in-stream PFOS signal into a distributed land-derived background that explains the upstream gradient and a residual at the lowest reaches that localizes the concentrated Wolverine source---an attribution that a grab sample, a probability survey, or a statistical occurrence model cannot make from the same data, because each integrates over the very pathways the process model separates. This is the sense in which the contribution is one of capability rather than goodness-of-fit: the process model supplies mechanism, mass balance, and---once a source is represented explicitly---counterfactuals, the quantities on which CERCLA allocation, Total Maximum Daily Load development, and intervention design actually operate. We are careful not to overstate the present fit: a log-space root-mean-square error of $0.15$~dex along the source-bearing mainstem ($0.74$~dex over all $20$ gauged reaches) on a single watershed is a first demonstration, not a validated predictive system, and the model's value at this stage is the auditable causal structure rather than the precision of any single predicted concentration. The uncertainty ensemble sharpens rather than softens this reading: that the predictive envelope brackets the mainstem ($P\text{-factor}=0.71$) but not the tributaries shows the tributary misfit to be structural---local soil and source heterogeneity---rather than a matter of unconstrained parameters.

\subsection{Why the engine's correctness checks matter}
PFAS fate-and-transport codes are difficult to verify because field data are sparse and confounded, so a model can fit observations for the wrong reasons. By anchoring the implementation to two reference checks that do not depend on field data---an independent high-precision solution of the soil equilibrium and numerical identity with the established SWAT\textsuperscript{+} channel solver---we move the burden of trust off the calibration and onto the code. The practical consequence is that when the calibrated model misfits, the misfit can be attributed to inputs (source representation, soil initialization) rather than to undiagnosed transport errors, as the structured Rogue residuals illustrate. This verification-first posture also makes the engine a reusable component: because the in-stream mathematics is provably the SWAT\textsuperscript{+} pesticide subset, the broad body of SWAT\textsuperscript{+} channel-process experience transfers directly.

\subsection{Transferability and the sediment-legacy frontier}
Because the engine is coupled to automated, national-data model generation, it is transferable to any U.S.\ watershed for which NHDPlus~HR and the national soil, land-use, and climate layers exist, with the per-watershed effort reduced to source representation and calibration rather than model construction. The most demanding future test is the post-source-control, sediment-legacy regime documented elsewhere in the Great Lakes basin, where water-column PFOS has fallen sharply after upstream control while fish-tissue concentrations have plateaued above advisory thresholds---a divergence attributed to contaminated sediment acting as a slowly releasing internal source. The engine already carries the necessary in-stream structure: an explicit benthic compartment with settling, resuspension, diffusion, and burial, evolving in time. Extending the demonstration to that regime requires adding a bioaccumulation tier and a measured sediment inventory, both beyond the present scope; we note it because the engine is positioned to test the sediment-legacy mechanism quantitatively, which equilibrium or source-inventory accounting cannot.

\subsection{Limitations}
Four limitations bound the claims. First, this version simulates PFOS only; the regulatory $4$~ng~L\textsuperscript{-1} MCL is PFOA/PFOS-specific and we make no class-wide claim. Second, the soil-PFAS initialization uses land-use and soil-class background ranges rather than measured per-site soil concentrations, which is what makes it national and automatable but also what produces the tributary scatter; a class-resolved or data-assimilated initialization is the clear next step. Third, the calibration is spatial, against per-station aggregate concentrations, not a temporal validation against a continuous concentration series, which the ambient record does not provide. Fourth, equifinality is managed---through global, class-based parameters and global sensitivity screening---not eliminated, and reported skill should be read as a workflow diagnostic on one watershed pending broader, multi-watershed and multi-analyte evaluation. Expert review remains appropriate before any decision use.
